%%%%%%%%%%%%%%%%%%%%%%%%%%%%%%%%%%%%%%%%%
% Proceedings of the National Academy of Sciences (PNAS)
% LaTeX Template
% Version 1.0 (19/5/13)
%
% This template has been downloaded from:
% http://www.LaTeXTemplates.com
%
% Original author:
% The PNAStwo class was created and is owned by PNAS:
% http://www.pnas.org/site/authors/LaTex.xhtml
% This template has been modified from the blank PNAS template to include
% examples of how to insert content and drastically change commenting. The
% structural integrity is maintained as in the original blank template.
%
% Original header:
%% PNAStmpl.tex
%% Template file to use for PNAS articles prepared in LaTeX
%% Version: Apr 14, 2008
%
%%%%%%%%%%%%%%%%%%%%%%%%%%%%%%%%%%%%%%%%%

%----------------------------------------------------------------------------------------
%	PACKAGES AND OTHER DOCUMENT CONFIGURATIONS
%----------------------------------------------------------------------------------------

%------------------------------------------------
% BASIC CLASS FILE
%------------------------------------------------

%% PNAStwo for two column articles is called by default.
%% Uncomment PNASone for single column articles. One column class
%% and style files are available upon request from pnas@nas.edu.

%\documentclass{pnasone}
\documentclass{pnastwo}

%------------------------------------------------
% POSITION OF TEXT
%------------------------------------------------

%% Changing position of text on physical page:
%% Since not all printers position
%% the printed page in the same place on the physical page,
%% you can change the position yourself here, if you need to:

% \advance\voffset -.5in % Minus dimension will raise the printed page on the 
                         %  physical page; positive dimension will lower it.

%% You may set the dimension to the size that you need.

%------------------------------------------------
% GRAPHICS STYLE FILE
%------------------------------------------------

%% Requires graphics style file (graphicx.sty), used for inserting
%% .eps/image files into LaTeX articles.
%% Note that inclusion of .eps files is for your reference only;
%% when submitting to PNAS please submit figures separately.

%% Type into the square brackets the name of the driver program 
%% that you are using. If you don't know, try dvips, which is the
%% most common PC driver, or textures for the Mac. These are the options:

% [dvips], [xdvi], [dvipdf], [dvipdfm], [dvipdfmx], [pdftex], [dvipsone],
% [dviwindo], [emtex], [dviwin], [pctexps], [pctexwin], [pctexhp], [pctex32],
% [truetex], [tcidvi], [vtex], [oztex], [textures], [xetex]

\usepackage{graphicx}

%------------------------------------------------
% OPTIONAL POSTSCRIPT FONT FILES
%------------------------------------------------

%% PostScript font files: You may need to edit the PNASoneF.sty
%% or PNAStwoF.sty file to make the font names match those on your system. 
%% Alternatively, you can leave the font style file commands commented out
%% and typeset your article using the default Computer Modern 
%% fonts (recommended). If accepted, your article will be typeset
%% at PNAS using PostScript fonts.

% Choose PNASoneF for one column; PNAStwoF for two column:
%\usepackage{PNASoneF}
%\usepackage{PNAStwoF}

%------------------------------------------------
% ADDITIONAL OPTIONAL STYLE FILES
%------------------------------------------------

%% The AMS math files are commonly used to gain access to useful features
%% like extended math fonts and math commands.

\usepackage{amssymb,amsfonts,amsmath}

%------------------------------------------------
% OPTIONAL MACRO FILES
%------------------------------------------------

%% Insert self-defined macros here.
%% \newcommand definitions are recommended; \def definitions are supported

%\newcommand{\mfrac}[2]{\frac{\displaystyle #1}{\displaystyle #2}}
%\def\s{\sigma}

%------------------------------------------------
% DO NOT EDIT THIS SECTION
%------------------------------------------------

%% For PNAS Only:
%\contributor{Submitted to Proceedings of the National Academy of Sciences of the United States of America}
%\url{www.pnas.org/cgi/doi/10.1073/pnas.0709640104}
%\copyrightyear{2008}
%\issuedate{Issue Date}
%\volume{Volume}
%\issuenumber{Issue Number}

%------------------------------------------------
% Languages problems (added by me)
%------------------------------------------------


\usepackage[french]{babel}
\usepackage[utf8]{inputenc}

%----------------------------------------------------------------------------------------

\begin{document}

%----------------------------------------------------------------------------------------
%	TITLE AND AUTHORS
%----------------------------------------------------------------------------------------

\title{Transcription factor binding sites detection in \textit{Paramecium}} % For titles, only capitalize the first letter

%------------------------------------------------

%% Enter authors via the \author command.  
%% Use \affil to define affiliations.
%% (Leave no spaces between author name and \affil command)

%% Note that the \thanks{} command has been disabled in favor of
%% a generic, reserved space for PNAS publication footnotes.

%% \author{<author name>
%% \affil{<number>}{<Institution>}} One number for each institution.
%% The same number should be used for authors that
%% are affiliated with the same institution, after the first time
%% only the number is needed, ie, \affil{number}{text}, \affil{number}{}
%% Then, before last author ...
%% \and
%% \author{<author name>
%% \affil{<number>}{}}

%% For example, assuming Garcia and Sonnery are both affiliated with
%% Universidad de Murcia:
%% \author{Roberta Graff\affil{1}{University of Cambridge, Cambridge,
%% United Kingdom},
%% Javier de Ruiz Garcia\affil{2}{Universidad de Murcia, Bioquimica y Biologia
%% Molecular, Murcia, Spain}, \and Franklin Sonnery\affil{2}{}}

\author{Matthias Grenié\affil{1}{ENS de Lyon, Département de Biologie},
Jean-François Goût\affil{2}{Indiana University, Biology Department, Lynch Lab}
\and
Michael Lynch\affil{2}{}}

\contributor{Submitted to Proceedings of the National Academy of Sciences
of the United States of America}

%----------------------------------------------------------------------------------------

\maketitle % The \maketitle command is necessary to build the title page

\begin{article}

%----------------------------------------------------------------------------------------
%	ABSTRACT, KEYWORDS AND ABBREVIATIONS
%----------------------------------------------------------------------------------------

\begin{abstract}
Lorem ipsum dolor sit amet, consectetur adipiscing elit. Sed vehicula metus sapien. Suspendisse pulvinar, felis ut hendrerit aliquet, dui nisi bibendum erat, fermentum mattis enim nibh id arcu. Vestibulum ultrices eros sed odio tincidunt bibendum. Pellentesque fermentum ante vel nulla commodo fermentum. Vestibulum in augue sit amet libero viverra accumsan eu at magna. Sed at ligula quis nibh pharetra facilisis non eu libero. Suspendisse non quam sit amet massa luctus interdum sit amet in purus. Integer id orci elit, vitae sollicitudin lectus.
\end{abstract}

%------------------------------------------------

\keywords{Keyword1 | Keyword2 | Keyword3} % When adding keywords, separate each term with a straight line: |

%------------------------------------------------

%% Optional for entering abbreviations, separate the abbreviation from
%% its definition with a comma, separate each pair with a semicolon:
%% for example:
%% \abbreviations{SAM, self-assembled monolayer; OTS,
%% octadecyltrichlorosilane}

% \abbreviations{}
\abbreviations{SAM, self-assembled monolayer; OTS, octadecyltrichlorosilane}

%----------------------------------------------------------------------------------------
%	PUBLICATION CONTENT
%----------------------------------------------------------------------------------------

%% The first letter of the article should be drop cap: \dropcap{} e.g.,
%\dropcap{I}n this article we study the evolution of ''almost-sharp'' fronts

\section*{Introduction}

Structure of the introduction
\newline
\begin{itemize}
\item Whole Genome Duplications background, major evolutionary force
\item the \textit{Paramecium} project, why \textit{Paramecium} is interesting, the aurelia complex
\item Here, focusing on the computational part, developing pipeline, showed that etc.
\end{itemize}

Since Ohno first hypothesized the influence of Whole-Genome Duplications (WGD), scientists kept showing that a broad number of organisms experienced several rounds of WGD: yeast, insects, Angiosperma, Vertebrates, Salmonids, and many others. WGD are evolutionary event when the genome of a given individual is duplicated, meaning that the whole genome is in two copies, duplicated pairs of genes are \textbf{paralogs}. WGD can also occur when two closely related species hybridize and form a fertile descendant.

WGDs may be involved in many evolutive radiations as they provide the raw material to explore new evolutionary landscapes. A recent study on the horshoe crab genome underlined that WGDs may be a more common phenomenon than stated. They showed that at least a WGD occurred in this conserved lineage, thus WGDs may not be evolutionary drivers.

Still, it seems that WGDs are widespread along the tree of life and since Ohno numerous models have been developed to explain the retention rate we observe between duplicate genes (reviewed in ). Gene Balance Hypothesis is for example a well explored hypothesis in the litterature, according to this hypothesis dosage-sensitive gene are more retained because of stœchiometry issues, it would thus explain the over retention of transcription factors and multicomplexes proteins.

To understand the consequences of WGDs we have been studying various \textit{Paramecium} species. \textit{Paramecium} are a Ciliates group (See position on phylogenetic tree). Aury \textit{et al.} showed that at least three round of WGDs occured in the \textit{Paramecium} genus, two of which occured in the \textit{aurelia} complex.  In the \textit{Paramecium} genus, the \textit{aurelia} complex of species is of particular interest. Indeed, 

The diversity of the \textit{Paramecium} ciliates is well studied. We focused on four species of \textit{Paramecium}: \textit{P. biaurelia}, \textit{P. sexaurelia}, \textit{P. tetraurelia} and \textit{P. caudatum} as an outgroup (see phylogenetic tree). The three \textit{aurelia} species underwent two rounds of WGDs, WGDX (... years ago) and WGDX (... years ago).

We have shown previous the tight link between gene expression and duplicate retention among \textit{P. caudatum} and \textit{P. tetraurelia}. Highly expressed genes are more retained than low expressed genes; the regulation of gene expression thus impacts gene retention. \textit{Cis}-regulatory sequences are upstream sequences influencing downstream gene expression.

Transcription Factor Binding Sites, regulatory sequences, regulatory network?, conserved sequences, studied for a long time. Regulatory sequences <-> WGD?

Objective: detect conserved TFBS among various our species -> to identify candidates for further studies.

\textbf{Biological questions:} Do gene expression is linked, in \textit{P.}, with specific motifs? How are TFBS affected by WGD? Are they conserved among species, is this linked to expression level? Conserved among each species? Is there a bias of TFBS usage in certain species?

%------------------------------------------------

\section*{Methods}

Used TranslatorX, PhyML, BigFoot, MEME.

We set up a pipeline to make our analyse (Fig. ), the code is available at.

\subsection*{Genomes and Annotation}

We used annotation and sequence from our previous analyses for P. and P. etc. (see .) and the additionnal sequence of.

\subsection*{Gene families}

Looking at the phylogenetic tree of the \textit{aurelia} species, two WGDs occured at the root and affected three of our species. We have established some gene families from WGD2 using comparison, each family contains a set of orthologous genes between the four \textit{aurelia} species studied, and eventually, the paralogous gene found in each species ; at maximum the families contain 13 different genes. Those families were established previously in our team. For details in the method see

\subsection*{Upstream sequences extraction}

We considered only families with at least 4 genes. Then, using our genome assembly and annotation of each species, we extracted upstream regions, upstream of the start codon, from 15nt with a cut-off at 250nt of all genes of the family. If the upstream region of a gene overlapped with another gene we discarded the gene, if the upstream region was less than 15nt long we also discarded the gene. Considering discarded genes, we kept only families with at least 4 members in our datasets.

\subsection*{Coding Sequences extraction and alignment}

Phylogenetic footprinting requires a phylogenetic tree to weigh the evolutionary signal of given motifs. A motif conserved between two close species will have less importance than a motif conserved in two distant species. Because we are focusing on the conservation of upstream sequences we chose not to use them to avoid circularity. Instead,corresponding coding sequences (CDS) were extracted and used to model phylogenetic trees for each family. We preferred to have a gene tree over species tree, to avoid eventual inconsistencies because of gene conversion (ref. needed). See Challenges section for explanations on the use of gene tree over species tree.

CDSs in each family were aligned using TranslatorX (ref. needed) a protein-guided alignment software. The Maximum Likelihood (ML) tree was then computed using PhyML (ref. needed) default parameters.

\subsection*{Phylogenetic footprinting}

We used a phylogenetic footprinting software BigFoot (ref. needed) to detect highly conserved motifs in upstream sequences. We used 10000 burn-in cycles and 20000 cycles with a sampling rate of 1000 for the Hidden Monte-Carlo Markov Chain (HMCMC) process. BigFoot aligns the given sequences with gaps and tries to identify conserved and non-conserved regions ; it models the evolution of those regions along the phylogenetic tree assuming conserved regions evolve more slowly than non-conserved ones. At the end of the analysis BigFoot outputs an alignment of sequences used to identify slow and fast evolving regions as well as, for each nucleotide in the alignment, the posterior probability of the alignment, higher values show higher confidence in the alignment, and the phylogenetic footprinting result, higher values indicating higher posterior probability of purifying selection.

Using a phylogenetic footprinting program means we have to use a phylogenetic tree and depending on the phylogenetic tree we are using, the evolutionary signal used in the footprinting is not identical.

The species tree (see figure.) gives us the relation between all considered, accounting for the various splits between species along with WGDs. The problem is that, not all gene families follow this tree. Because of the successive round of WGDs there are several fates possible for pair of duplicate after the first WGD. Some of these genes may cluster together in the same leaves, if the pairs diverge between species ; another possible outcome is the subfunctionnalization of each gene before the second WGD, meaning before the speciation of the \textit{aurelia} complex, thus genes from different species would cluster together in a phylogenetic tree ; or even a combination of the previous outcome and gene conversion, leading paralogs to recombine in a copy-paste way, changing dramatically the gene tree.

Using these scores we detected motifs of at least 6 nucleotides long, alignment score over 0.8 and phylogenetic score over 0.9. Because of the known biological nature of Transcription Factor Binding Sites, we allowed for a "gap" in motifs of 2 nucleotide, so that the scores could drop under the thresholds in those gaps.

BigFoot does not output directly identified motifs, instead it produces two files with an alignment of the sequences and associated phylogenetic and alignment scores, as explained above.

Transcription factor binding sites are known to be generally conserved but degenerate on certain positions. For example, ... showed that this motif was conserved ....NN... with two highly variable positions (denoted by "N", meaning "A", "T", "C" or "G" using IUPAC notation). Thus, to seek biologically relevant motifs, we had to take into account that in the middle of motifs, the phylogenetic score could drop on several positions, before rising again.

To answer this problem we use a sliding window method of 8 nucleotides in our analysis: for each family, we looked at the scores of 8 nucleotides at the time and slide along the sequences. If the window contained at least 6 bases with scores above our thresholds, we would retain this motif. Then, from this particular region we would try to extend the sequence by adding adjacent nucleotides with good scores.

\subsection*{Comparison with MEME}

To check our predictions and assess the conservation of found motifs, we compared motifs prediction with those of MEME, a widely used \textit{ab initio} motif finding tools. It searches for statistically significant motifs, with a gap-less, local multi-alignment system.

MEME was shown to have a very high False Positive Rate of the discovery (Ref. needed). That is why many studies combine multiple motif detection tools (Ref. needed). In our case, MEME is of particular interest as it obtains motifs using a totally different method from BigFoot.

To assess the relevance of motifs found from BigFoot's outputs, we compared them to a well-known motif finding program: MEME. For each family we identified overlapping motifs between MEME and BigFoot. We computed an overlapping index as follows: 
\begin{equation}
0 \leq \frac{nucleotides in common}{size of the smallest motif} \leq 1
\end{equation}
if this index was over $0.9$ we would then considered the motif as relevant.

\subsection*{Motif classification and data analysis}

All the analyses were produced using R, scatter plots and graphs were produced using the R package ggplot2.



%------------------------------------------------

\section*{Pipeline}

Description of the full pipeline

%------------------------------------------------

\section{Challenges}

\subsection{Species tree or gene tree?}

Using a phylogenetic footprinting program means we have to use a phylogenetic tree and depending on the phylogenetic tree we are using, the evolutionary signal used in the footprinting is not identical.

The species tree (see figure.) gives us the relation between all considered, accounting for the various splits between species along with WGDs. The problem is that, not all gene families follow this tree. In particular, gene conversion, is heavily involved in shifting genes tree away from species tree: paralogs gene recombine and ...

\subsection{Motif detection strategy}

BigFoot does not output directly identified motifs, instead it produces two files with an alignment of the sequences and associated phylogenetic and alignment scores. The phylogenetic scores is computed, etc. The alignement scores, etc.

Transcription factor binding sites are known to be generally conserved but degenerate on certain positions. For example, ... showed that this motif was conserved ....NN... with two highly variable positions (denoted by\"N\", meaning \"A\", \"T\", \"C\" or \"G\" using IUPAC notation). Thus, to seek biologically relevant motifs, we had to take into account that in the middle of motifs, the phylogenetic score could drop on several positions. (Show a phylogenetic score profile?)

To answer this problem we use a sliding window method of 8 nucleotides in our analysis: for each family, we looked at the scores of 8 nucletodies at the time and slide along the sequences. If the window contained at least 6 bases with scores above our thresholds, we would retain this motif. Then, from this particular region we would try to extend the sequence by adding adjacent nucleotides with good scores.

\subsection{Measure motif relevance}

To assess the relevance of motifs found from BigFoot's outputs, we compared them to a well-known motif finding program: MEME

%------------------------------------------------

\section*{Results}
\label{sec:Results}

From our orthologs and paralogs gene familes we had 5751 families, only 5008 gene families had at least 4 genes. We ran our pipeline on these 5008 families, to detect conserved motifs among upstream sequences of \textit{Paramecium} genomes. The pipeline was set to detect motifs, in upstream regions ranging between 15nt and 250nt (see~\nameref{sec:Methods}) , of at least 6 nucleotides with an eventual gap in conservation (see ~\autoref{fig:DegenerateMotif}) for more biological relevance.

Using BigFoot we identified 811 different motifs in 608 families. Comparing these motifs with MEME ones for partial match, only 117 unique motifs were retained. For each of these, we extracted the name of all genes containing the motif, in each species. To focus only on conserved motifs, we conserved only motifs matching in all the four species. Only 10 motifs matched these conditions. We then computed the ratio of genes containing the motif over the total number of genes (results are summed up in~\autoref{tab:Motifs}).

%------------------------------------------------

\section*{Future Directions}
\label{sec:Future}

We identified motifs using an automated pipeline (Fig.~\ref{fig:Pipeline}), comparing motifs identified using phylogenetic footprinting and statistical enrichment methods. We were able to dectect X motifs, associated with specific expression values and distance. Still, we plan to improve this pipeline in several ways.

After identifying motifs, we scan the entire \textit{Paramecium} genomes to find associated genes, however we only search for perfect matches. Thus, we do not take into account degenerate motifs. Thus we can miss genes that have slightly different motifs. Recently XXX showed that looking exact matches during motif identification may miss most of the signal. Indeed we have increase evidence that real motifs are most of the time degenerate and that Transcription Factor binding mostly depends on the general DNA context.

After identifying motifs, one's might want to group the motifs by similarity. We generally think about motifs as families, having specific sequences with some variation. For the moment, in our pipeline, we do not cluster motifs using "supposed" family, we consider each exact motifs as unique. However, several clustering methods exist to measure the diversity among motifs. The small sizes of motifs (about 6nt to 15nt) make these methods complicated to develop. Several methods have been implemented, but still, are not adapted for motifs and do not output families of motifs. 

Another way to think about the problem can be to cluster motifs according to features one can measure on them. For example, one can compute several statistical indexes on motifs to compare. Thus we do not need to implement a new algorithm to cluster motifs as other classical methods such as principal components analysis, hierarchal clustering or \textit{k}-means clustering can be applied. Still, these methods rely on the indexes you choose to measure, the motifs you are studying have to have well-defined clusters using those variables. We tried to cluster the studied motifs with this hierarchal clustering method (Fig.).

Several other new methods have been proposed to cluster motifs as they are, considering the sequences and the variability along each nucleotide. Some of them rely on position-specific weight matrices (PSM) which are matrices of the abundances of each base at each position. In a future improvement of the pipeline we cool try to implement such methods to reduce the redundancy of our data.

It has been proven that most motif identification tools have a high false positive discovery rate, thus, we need ways to confirm the relevance of found motifs. One way used by \citealt{liseron-monfils_promzea:_2013} is to increase the number of tools used in the pipeline, in their analysis they used MEME, BioProspector and Weeder and filtered out significant results and then only combined them. They have selected in each set, significant motifs using p-values and then only compare them. Liseron \textit{et al.} showed that by combining these tools they were able to increase the sensitivity by 22\% over the best standalone tool. Here we used MEME and BigFoot but we could include other tools in our workflow.
%----------------------------------------------------------------------------------------
%	MATERIALS AND METHODS
%----------------------------------------------------------------------------------------

%% Optional Materials and Methods Section
%% The Materials and Methods section header will be added automatically.


%----------------------------------------------------------------------------------------
%	ACKNOWLEDGEMENTS
%----------------------------------------------------------------------------------------

\begin{acknowledgments}
Thanks to M. Lynch and J.-F. Goût. Funding by NSF.
\end{acknowledgments}

%----------------------------------------------------------------------------------------
%	BIBLIOGRAPHY
%----------------------------------------------------------------------------------------

%% PNAS does not support submission of supporting .tex files such as BibTeX.
%% Instead all references must be included in the article .tex document. 
%% If you currently use BibTeX, your bibliography is formed because the 
%% command \verb+\bibliography{}+ brings the <filename>.bbl file into your
%% .tex document. To conform to PNAS requirements, copy the reference listings
%% from your .bbl file and add them to the article .tex file, using the
%% bibliography environment described above.  

%%  Contact pnas@nas.edu if you need assistance with your
%%  bibliography.

% Sample bibliography item in PNAS format:
%% \bibitem{in-text reference} comma-separated author names up to 5,
%% for more than 5 authors use first author last name et al. (year published)
%% article title  {\it Journal Name} volume #: start page-end page.
%% ie,
% \bibitem{Neuhaus} Neuhaus J-M, Sitcher L, Meins F, Jr, Boller T (1991) 
% A short C-terminal sequence is necessary and sufficient for the
% targeting of chitinases to the plant vacuole. 
% {\it Proc Natl Acad Sci USA} 88:10362-10366.


%% Enter the largest bibliography number in the facing curly brackets
%% following \begin{thebibliography}

\begin{thebibliography}{10}
\bibitem{BN}
M.~Belkin and P.~Niyogi, {\em Using manifold structure for partially
  labelled classification}, Advances in NIPS, 15 (2003).

\bibitem{BBG:EmbeddingRiemannianManifoldHeatKernel}
P.~B\'erard, G.~Besson, and S.~Gallot, {\em Embedding {R}iemannian
  manifolds by their heat kernel}, Geom. and Fun. Anal., 4 (1994),
  pp.~374--398.

\bibitem{CLAcha1}
R.R.~Coifman and S.~Lafon, {\em Diffusion maps}, Appl. Comp. Harm. Anal.,
  21 (2006), pp.~5--30.

\bibitem{DiffusionPNAS}
R.R.~Coifman, S.~Lafon, A.~Lee, M.~Maggioni, B.~Nadler, F.~Warner, and
  S.~Zucker, {\em Geometric diffusions as a tool for harmonic analysis and
  structure definition of data. {P}art {I}: Diffusion maps}, Proc. of Nat.
  Acad. Sci.,  (2005), pp.~7426--7431.

\bibitem{Clementi:LowDimensionaFreeEnergyLandscapesProteinFolding}
P.~Das, M.~Moll, H.~Stamati, L.~Kavraki, and C.~Clementi, {\em
  Low-dimensional, free-energy landscapes of protein-folding reactions by
  nonlinear dimensionality reduction}, P.N.A.S., 103 (2006), pp.~9885--9890.

\bibitem{DoGri}
D.~Donoho and C.~Grimes, {\em Hessian eigenmaps: new locally linear
  embedding techniques for high-dimensional data}, Proceedings of the National
  Academy of Sciences, 100 (2003), pp.~5591--5596.

\bibitem{DoGri:WhenDoesIsoMap}
D.~L. Donoho and C.~Grimes, {\em When does isomap recover natural
  parameterization of families of articulated images?}, Tech. Report Tech. Rep.
  2002-27, Department of Statistics, Stanford University, August 2002.

\bibitem{GruterWidman:GreenFunction}
M.~Gr\"uter and K.-O. Widman, {\em The {G}reen function for uniformly
  elliptic equations}, Man. Math., 37 (1982), pp.~303--342.

\bibitem{Simon:NeumannEssentialSpectrum}
R.~Hempel, L.~Seco, and B.~Simon, {\em The essential spectrum of neumann
  laplacians on some bounded singular domains}, 1991.

\bibitem{1}
Kadison, R.\ V.\ and Singer, I.\ M.\ (1959)
Extensions of pure states, {\it Amer.\ J.\ Math.\ \bf
81}, 383-400.

\bibitem{2}
Anderson, J.\ (1981) A conjecture concerning the pure states of
$B(H)$ and a related theorem. in {\it Topics in Modern Operator
Theory}, Birkha\"user, pp.\ 27-43.

\bibitem{3}
Anderson, J.\ (1979) Extreme points in sets of
positive linear maps on $B(H)$. {\it J.\ Funct.\
Anal.\
\bf 31}, 195-217.

\bibitem{4}
Anderson, J.\ (1979) Pathology in the Calkin algebra. {\it J.\
Operator Theory \bf 2}, 159-167.

\bibitem{5}
Johnson, B.\ E.\ and Parrott, S.\ K.\ (1972) Operators commuting
with a von Neumann algebra modulo the set of compact operators.
{\it J.\ Funct.\ Anal.\ \bf 11}, 39-61.

\bibitem{6}
Akemann, C.\ and Weaver, N.\ (2004) Consistency of a
counterexample to Naimark's problem. {\it Proc.\ Nat.\ Acad.\
Sci.\ USA \bf 101}, 7522-7525.

\bibitem{TSL}
J.~Tenenbaum, V.~de~Silva, and J.~Langford, {\em A global geometric
  framework for nonlinear dimensionality reduction}, Science, 290 (2000),
  pp.~2319--2323.

\bibitem{ZhaZha}
Z.~Zhang and H.~Zha, {\em Principal manifolds and nonlinear dimension
  reduction via local tangent space alignement}, Tech. Report CSE-02-019,
  Department of computer science and engineering, Pennsylvania State
  University, 2002.
\end{thebibliography}

%----------------------------------------------------------------------------------------

\end{article}

%----------------------------------------------------------------------------------------
%	FIGURES AND TABLES
%----------------------------------------------------------------------------------------

%% Adding Figure and Table References
%% Be sure to add figures and tables after \end{article}
%% and before \end{document}

%% For figures, put the caption below the illustration.
%%
%% \begin{figure}
%% \caption{Almost Sharp Front}\label{afoto}
%% \end{figure}

\begin{figure}[h]
\centerline{\includegraphics[width=0.4\linewidth]{placeholder.jpg}}
\caption{Figure caption}\label{placeholder}
\end{figure}

%% For Tables, put caption above table
%%
%% Table caption should start with a capital letter, continue with lower case
%% and not have a period at the end
%% Using @{\vrule height ?? depth ?? width0pt} in the tabular preamble will
%% keep that much space between every line in the table.

%% \begin{table}
%% \caption{Repeat length of longer allele by age of onset class}
%% \begin{tabular}{@{\vrule height 10.5pt depth4pt  width0pt}lrcccc}
%% table text
%% \end{tabular}
%% \end{table}

\begin{table}[h]
\caption{Table caption}\label{sampletable}
\begin{tabular}{l l l}
\hline
\textbf{Treatments} & \textbf{Response 1} & \textbf{Response 2}\\
\hline
Treatment 1 & 0.0003262 & 0.562 \\
Treatment 2 & 0.0015681 & 0.910 \\
Treatment 3 & 0.0009271 & 0.296 \\
\hline
\end{tabular}
\end{table}

%% For two column figures and tables, use the following:

%% \begin{figure*}
%% \caption{Almost Sharp Front}\label{afoto}
%% \end{figure*}

%% \begin{table*}
%% \caption{Repeat length of longer allele by age of onset class}
%% \begin{tabular}{ccc}
%% table text
%% \end{tabular}
%% \end{table*}

%----------------------------------------------------------------------------------------

\end{document}